% Generated from tex/chapters/ch04/07_実務上の考慮.tex for the SWEBOK structure tree
\subsubsection{構築品質}

構築品質は、コードやその近くにある成果物の品質に焦点を当てる。要求や上位設計に関する品質活動も重要だが、構築品質では、詳細設計、コード、単体テスト、統合テスト、ビルド、静的解析、レビューが中心になる。

構築品質を高める主な技法には、単体テスト、統合テスト、テストファースト開発、表明、防御的プログラミング、デバッグ、コードインスペクション、技術レビュー、セキュリティレビュー、静的解析がある。

特にセキュリティでは、よく知られた脆弱性パターンを避ける必要がある。入力検証不足、境界チェック不足、例外の握りつぶし、危険なライブラリの利用、認可確認漏れなどは、構築段階で入りやすい問題である。
